%! Logarithmic Function Manipulation — Unit 5, Lesson 5.4 — Warm-Up
\documentclass[10pt]{article}
\usepackage{precalculus-article}
\usepackage{precalculus-boxes}
\newcommand{\TallMath}[1]{$\displaystyle #1\rule[-1.4em]{0pt}{3.2em}$}

\begin{document}

\pageheader{Unit 5, Lesson 5.4}{Warm-Up}
\namedateperiod

\vspace{0.15in}

\begin{enumerate}[leftmargin=1.8em, itemsep=14pt]

\item Evaluate each logarithm.

\vspace{6pt}
\begin{center}
\renewcommand{\arraystretch}{2.0}
\begin{tabular}{@{} c @{\qquad} c @{\qquad} c @{}}
  $\log_2(32) = $ \blank{2cm} &
  $\log_3(81) = $ \blank{2cm} &
  $\log_{10}(1000) = $ \blank{2cm} \\
\end{tabular}
\end{center}

\item Let $f(x) = \log_5(x)$. Evaluate $f(25)$, $f(5)$, and $f(1)$.

\vspace{6pt}
\begin{center}
\renewcommand{\arraystretch}{2.0}
\begin{tabular}{@{} c @{\qquad} c @{\qquad} c @{}}
  $f(25) = $ \blank{2cm} &
  $f(5) = $ \blank{2cm} &
  $f(1) = $ \blank{2cm} \\
\end{tabular}
\end{center}

\vspace{4pt}
What pattern do you notice in the outputs as the input multiplies by 5?

\writeline

\item Recall the exponent rule: $2^3 \cdot 2^4 = 2^{\blank{1cm}}$.

\vspace{6pt}
Based on this, predict: $\log_2(8 \cdot 16) = \log_2(8) + \log_2(16) = \blank{1.5cm} + \blank{1.5cm} = \blank{1.5cm}$.

\vspace{4pt}
Verify by evaluating $\log_2(128)$ directly: $\log_2(128) = $ \blank{2cm}.

\end{enumerate}

\end{document}
